\chapter{Images, textures and sprites}

\section{Two things that sound identical and are not}

\begin{center}
\begin{tabular}{@{}lll@{}}
\toprule
 & \textbf{\lstinline|Image|} & \textbf{\lstinline|Texture2D|} \\
\midrule
Lives in & normal RAM (CPU) & video memory (GPU) \\
You can  & read and edit pixels & draw it, fast \\
You cannot & draw it & edit single pixels \\
Typical use & load, resize, recolour & everything you see \\
\bottomrule
\end{tabular}
\end{center}

The usual flow is: get an \lstinline|Image|, upload it once, throw the
\lstinline|Image| away, draw the \lstinline|Texture2D| forever.

\begin{lstlisting}
Image img = LoadImage("player.png");        // into RAM
ImageResize(&img, 64, 64);                  // cheap to edit here
Texture2D tex = LoadTextureFromImage(img);  // upload to the GPU
UnloadImage(img);                           // the RAM copy is dead weight now
\end{lstlisting}

If you do not need to touch the pixels, \lstinline|LoadTexture("player.png")|
does all four steps for you.

\section{The golden rule of loading}

\gotcha{Load \textbf{outside} the loop, unload \textbf{after} the loop. Calling
\lstinline|LoadTexture()| inside the frame loop uploads a fresh copy to the GPU
sixty times a second. A minute later you are out of video memory and you have no
idea why.}

\begin{lstlisting}
Texture2D tex = LoadTexture("player.png");   // once

while (!WindowShouldClose())
{
    BeginDrawing();
        DrawTexture(tex, 100, 100, WHITE);   // every frame: just draws
    EndDrawing();
}

UnloadTexture(tex);                          // once
CloseWindow();
\end{lstlisting}

The same rule covers fonts, sounds, music and models. Every
\lstinline|Load...| has a matching \lstinline|Unload...|, and they belong on
either side of the loop.

\section{Drawing a texture, four ways}

\begin{lstlisting}
DrawTexture(tex, x, y, WHITE);                       // plain, at a position
DrawTextureV(tex, position, WHITE);                  // same, Vector2
DrawTextureEx(tex, position, rotation, scale, WHITE);// rotated and scaled
DrawTextureRec(tex, source, position, WHITE);        // just a PIECE of it
DrawTexturePro(tex, source, dest, origin, rot, WHITE); // everything at once
\end{lstlisting}

That \lstinline|WHITE| at the end is the \emph{tint}: the texture's colours get
multiplied by it. \lstinline|WHITE| means ``unchanged''. Pass \lstinline|RED| and
everything goes red; pass \lstinline|Fade(WHITE, 0.5f)| and the sprite is half
transparent. This is how you flash a character on damage without a second image.

\section{Sprite sheets}

An animation is not many files, it is one image holding many frames side by side,
plus a \lstinline|Rectangle| that says which frame to show:

\begin{center}
\begin{tikzpicture}[scale=0.9,every node/.style={font=\sffamily\footnotesize}]
  \foreach \i in {0,1,2,3,4,5} {
    \draw[codeframe,fill=rlight] (\i*1.5,0) rectangle (\i*1.5+1.5,1.5);
    \node[text=rgrey] at (\i*1.5+0.75,0.75) {\i};
  }
  \draw[rred,line width=1.4pt] (3,0) rectangle (4.5,1.5);
  \node[text=rred] at (3.75,-0.4) {source rectangle};
  \node[text=rgrey] at (4.5,1.95) {one texture, six frames};
\end{tikzpicture}
\end{center}

\begin{lstlisting}
Rectangle source = { frame*FRAME_SIZE, 0, FRAME_SIZE, FRAME_SIZE };
DrawTextureRec(sheet, source, position, WHITE);
\end{lstlisting}

And you advance \lstinline|frame| on a timer, never on a frame counter, for the
reason chapter 5 gave:

\begin{lstlisting}
frameTimer += dt;
while (frameTimer >= frameDuration)
{
    frameTimer -= frameDuration;
    frame = (frame + 1) % FRAME_COUNT;
}
\end{lstlisting}

\gotcha{To flip a sprite horizontally, make the source rectangle's
\lstinline|width| \emph{negative}. It reads like a bug and it is the intended
API: \lstinline|source.width *= -1| and your character faces the other way,
with no second image and no rotation.}

\section{Where files are looked for}

Paths are relative to the \emph{working directory}, which is wherever you
launched the program from, not where the executable sits. That is why running
from an IDE sometimes finds your \file{.png} and running from a terminal does
not.

This project ships a helper for exactly that, in \file{src/resource\_dir.h}:

\begin{lstlisting}
SearchAndSetResourceDir("resources");   // find it and chdir there
Texture2D tex = LoadTexture("wabbit_alpha.png");
\end{lstlisting}

\gotcha{A texture that failed to load is not a crash: you get a texture with
\lstinline|id == 0| and you draw nothing, silently. If a sprite is invisible,
check \lstinline|IsTextureValid(tex)| before suspecting your draw code.}

\tryit{Lesson 5 builds a six frame sprite sheet in memory --- no external files
--- and plays it back, showing the whole sheet, the current frame, and the source
rectangle moving across it.}{./course2d/build.sh 5}
